\subsection{কোড থেকে তথ্য উদ্ধার ও ধারণাভিত্তিক উদাহরণ}
সংখ্যা পদ্ধতি সংখ্যার মান প্রকাশ করে; কোড কোনো নির্ধারিত তথ্যের প্রতিনিধিত্ব করে। একই বিটগুলো কোন নিয়মে লেখা হয়েছে তা না জানলে সঠিক অর্থ উদ্ধার করা যায় না।

\chThreeExample{১}{BCD থেকে দশমিক এবং ভুল group}
$(0101\;1000\;0011)_{\mathrm{BCD}}$-এর দশমিক মান কত?
\chThreeAnswer চার বিট করে পড়লে অঙ্ক ৫, ৮, ৩; ফল $(583)_{10}$। পুরো ১২ বিটকে একটি সাধারণ বাইনারি সংখ্যা ধরে পড়া যাবে না। $1010$ থেকে $1111$ পর্যন্ত ছয়টি group 8421 BCD-তে কোনো দশমিক অঙ্কের বৈধ কোড নয়। তাই $0101\;1010$-কে বৈধ দুই অঙ্কের BCD বলা যাবে না।

\chThreeExample{২}{তিন রকম প্রকাশে একই ৩৯}
\chThreeAnswer
\[(39)_{10}=(100111)_2,\quad(39)_{\mathrm{BCD}}=0011\;1001.\]
BCD-তে ৮ বিট লাগে; বাইনারির সর্বনিম্ন দৈর্ঘ্য ৬ বিট। ASCII-তে দুই ক্যারেক্টার \textenglish{3} ও \textenglish{9}-এর কোড যথাক্রমে ৫১ ও ৫৭, আট বিটে $00110011\;00111001$। লেখা \textenglish{39}, সংখ্যামান ৩৯ এবং BCD প্রকাশ একই জিনিস নয়।

\chThreeExample{৩}{Excess-3 থেকে দশমিক}
$1000\;0111$-এর XS-3 মান উদ্ধার করো।
\chThreeAnswer group-এর বাইনারি মান ৮ ও ৭। প্রতিটি থেকে ৩ বাদ দিলে ৫ ও ৪; ফল ৫৪। পুরো সংখ্যার মান থেকে একবার ৩ বিয়োগ করা যাবে না।

\chThreeExample{৪}{Gray থেকে বাইনারি}
Gray কোড $1101$-এর বাইনারি মান বের করো।
\chThreeAnswer প্রথম বাইনারি বিট Gray-এর প্রথম বিট, ১। এরপর আগের পাওয়া \emph{বাইনারি} বিটের সঙ্গে পরের Gray বিট XOR করতে হয়।
\[b_3=1,\quad b_2=1\oplus1=0,\quad b_1=0\oplus0=0,\quad b_0=0\oplus1=1.\]
অতএব বাইনারি $1001$। ৭ থেকে ৮-এ সাধারণ বাইনারিতে $0111\to1000$-এ চার বিট বদলায়; reflected Gray-তে $0100\to1100$-এ এক বিট বদলায়। অবস্থান শনাক্তের সময় একাধিক বিট একসঙ্গে বদলানোর অস্পষ্টতা কমাতে এ বৈশিষ্ট্য কাজে লাগে।

\smallhead{ASCII, EBCDIC ও Unicode: প্রয়োজনীয় পার্থক্য}
মূল ASCII একটি ৭-বিট মানদণ্ড; ১২৮টি কোডে নিয়ন্ত্রণ চিহ্নও অন্তর্ভুক্ত। ৮-বিটের বিভিন্ন code page-কে সাধারণভাবে extended ASCII বলা হয়; এগুলো সব একই আন্তর্জাতিক সারণি নয়। EBCDIC-এর code page-ও ভিন্ন হতে পারে। \textenglish{Python}-এর \textenglish{cp037} EBCDIC code page-এ \textenglish{A} হলো হেক্সাডেসিমাল C1, আর ASCII-তে 41; কোডপদ্ধতি না জানলে অক্ষর ভুল পড়া যায়।

Unicode প্রতিটি ক্যারেক্টারের জন্য \textbf{কোড পয়েন্ট} নির্ধারণ করে। বাংলা অ-এর কোড পয়েন্ট \textenglish{U+0985}। এটি সংরক্ষণের জন্য UTF-8, UTF-16 বা UTF-32 ব্যবহার করা যায়। সব Unicode ক্যারেক্টার সবসময় ১৬ বিটে সংরক্ষিত হয়—এ দাবি সঠিক নয়। প্রাথমিক ১৬-বিট ধারণা অনুযায়ী $2^{16}=65536$টি মান বলা হয়; বর্তমান সম্পূর্ণ Unicode-কে ওই সীমায় বেঁধে বলা যাবে না।

\chThreeExample{৫}{বাংলা অ সংরক্ষণে কত বাইট?}
\chThreeAnswer UTF-8-এ অ-এর bytes হলো হেক্সাডেসিমাল \textenglish{E0 A6 85}, অর্থাৎ ৩ বাইট। UTF-16-এ এর জন্য একটি ১৬-বিট code unit, অর্থাৎ ২ বাইট লাগে। UTF-32-এ ৪ বাইট। কোড পয়েন্ট ও সংরক্ষণের বিটসংখ্যা পৃথক ধারণা।

\chThreeExample{৬}{ধারণাভিত্তিক MCQ}
তিন অঙ্কের দশমিক সংখ্যার 8421 BCD প্রকাশে কত বিট লাগে? (ক) ৩, (খ) ৮, (গ) ১০, (ঘ) ১২।
\chThreeAnswer (ঘ) ১২; প্রতিটি অঙ্কে ৪ বিট। তিন অঙ্কের সাধারণ বাইনারি প্রকাশের বিটসংখ্যা মানের ওপর নির্ভর করে।

% Mahbub teaching-type follow-up, 2026-10-14
\chThreeExample{৭}{ASCII-তে অক্ষর ও অঙ্কের লেখা}
\textenglish{B5} লেখাটির ASCII bytes লেখো। $01000010\;00110101$ bytes থেকে লেখাটি উদ্ধার করো।
\chThreeAnswer ASCII-তে B-এর মান ৬৬ এবং 5 ক্যারেক্টারের মান ৫৩। আট বিটে যথাক্রমে 01000010 ও 00110101, hexadecimal-এ 42 ও 35। বিপরীতভাবে প্রতিটি byte আলাদা ASCII মান হিসেবে পড়লে B5 পাওয়া যায়। সংখ্যামান ৫-এর binary 101; ক্যারেক্টার 5-এর ASCII 00110101।

\chThreeExample{৮}{BCD-তে শুরুতে শূন্য এবং সাধারণ বাইনারিতে রূপান্তর}
$0000\;0100\;0110$-কে 8421 BCD ধরে দশমিক ও সাধারণ বাইনারিতে লেখো।
\chThreeAnswer অঙ্কগুলো ০, ৪, ৬; লেখা 046, সংখ্যামান ৪৬। তাই সাধারণ বাইনারি $101110$। তিন অঙ্কের লেখা রাখতে BCD-তে ১২ বিট লাগে, কিন্তু সংখ্যামানের সর্বনিম্ন binary দৈর্ঘ্য ৬ বিট। একই ১২ বিটকে সরাসরি binary ধরে পড়লে $(046)_{16}=70$ হবে; তা BCD decoding নয়।
